\subsection{Local Brauer support at seventeen, after McCallum}
\label{subsec:seventeen-local-brauer}

This subsection proves a single-primary local condition using
McCallum's cusp-disc argument, then proves full nonabelian localization
over a regular-prime arithmetic source.  Write $F_p:X^p+Y^p=Z^p$; for odd $p$ this is isomorphic over
$\mathbb Q$ to $X_p$ by changing the sign of $Z$.  Let $\mathcal F_p$
be its proper plane model over $\mathbb Z$. The rational boundary
$D_p(\mathbb Q)$ consists of its three rational coordinate cusps.

For a proper curve $C/\mathbb Q$ with proper model $\mathcal C$, put
\[
 \operatorname{Br}_{\{p\}}(\mathcal C)
 =\{a\in\operatorname{Br}(C):
 a_v\text{ extends over }\mathcal C_{\mathbb Z_v}
 \text{ for every finite }v\ne p,\quad a_\infty=0\}.
\]
Its local zero set consists of the $x\in C(\mathbb Q_p)$ satisfying
$\operatorname{inv}_p a(x)=0$ for every such $a$.  Evaluation outside
$p$ vanishes by properness and $\operatorname{Br}(\mathbb Z_v)=0$.
Thus this zero set contains the projection of the Brauer--Manin set,
and every rational point belongs to it
\citep[\S1, pp.~393--394]{McCallum2001Brauer}.

\begin{proposition}[The regular-prime cusp residue classes]
\label{prop:regular-primary-brauer-cusp-discs}
Let $p>5$ be regular, meaning that the class number of
$\mathbb Q(\zeta_p)$ is prime to $p$.  In each residue class of a
rational coordinate cusp, the local zero set
$F_p(\mathbb Q_p)^{\operatorname{Br}_{\{p\}}(\mathcal F_p)[p^\infty]}$
contains just that cusp.  In particular,
\[
 F_{17}(\mathbb Q_{17})^{\operatorname{Br}_{\{17\}}(\mathcal F_{17})[17^\infty]}
       =D_{17}(\mathbb Q).
\]
\end{proposition}
\begin{proof}
We first explain the extension of McCallum's regular-prime argument
from rational points to this local zero set.  Put
$K=\mathbb Q(\zeta_p)$, let $\lambda$ be its unique prime over $p$,
and set $\pi=1-\zeta_p$ and $\Delta=\operatorname{Gal}(K/\mathbb Q)$.
Choose a wild cyclic quotient $C$ of $F_p$, with rational branch
cusps $O_0,O_1,O_\infty$, and write $A=\operatorname{Jac}(C)$.
Here wild means that, for branch exponents $a+b+c=0$,
$\gamma=a^ab^bc^c$ satisfies $\gamma^{p-1}\not\equiv1\pmod {p^2}$.
Such a quotient exists by
\citet[Lemma~9, pp.~580--581]{McCallum1994}.
Use $O_0$ as Abel--Jacobi base point and write $j:C\to A$.

The regularity hypothesis supplies classes
\[
 d_k\in H^1(K,A)[\pi^k]\qquad(k\ge1)
\]
which vanish locally away from $\lambda$, satisfy
$\pi^{k-1}d_k=d_1$, and whose first member pairs nontrivially with
$(A(K_\lambda)/\pi A(K_\lambda))^\Delta$.
These are exactly the classes of
\citet[Lemma~10 and Theorem~5, pp.~581--582]{McCallum1994}.
To specify where regularity is used, the construction begins with
the Kummer class of $\zeta_p$ and lifts it through
$H^1(\mathcal O_K[1/p],A[\pi^k])$.
The obstruction $H^2(\mathcal O_K[1/p],\mu_p)$ vanishes: its
Picard contribution vanishes by regularity, and its Brauer
contribution vanishes because the only allowed finite invariant
is at $\lambda$, all infinite places are complex, and the global
sum of invariants is zero.

Let $B$ be the common annihilator of these local classes in
$A(\mathbb Q_p)$:
\[
 B=\{P:\langle\operatorname{res}P,d_{k,\lambda}\rangle=0
                     \text{ for every }k\ge1\}.
\]
The classes $d_k$ are $p$-primary, because $p$ is a unit times
$\pi^{p-1}$. The canonical principal polarization of $A$, defined
 over $\Q$, identifies these with classes for the dual abelian variety.
The global classes $e_k=\operatorname{cor}_{K/\mathbb Q}d_k$
vanish at every place other than $p$.  Compatibility of the Tate
pairing with restriction and corestriction gives
\[
 \langle P,e_{k,p}\rangle_{\mathbb Q_p}
       =\langle\operatorname{res}P,d_{k,\lambda}\rangle_{K_\lambda}.
\]
There is no degree factor in this identity.  The relation between
the Tate pairing and Brauer evaluation, including the choice of
representatives with zero evaluations away from $p$, is
\citet[Lemma~1 and equation~(5), author pp.~4--5]{McCallum2001Brauer}.
Normalize the Brauer representative $a_k$ of $e_k$ by
$a_k(O_0)=0$. The rational base point splits the constants in
$\operatorname{Br}(C)$, and Hochschild--Serre identifies its quotient
by constants with $H^1(\Q,A)$. If $p^me_k=0$, then $p^ma_k$ is
constant and evaluates to zero at $O_0$, so $p^ma_k=0$. At every
$v\ne p$, the vanishing of $e_{k,v}$ makes $a_{k,v}$ constant;
normalization makes it the zero class. Thus these are genuinely
$p$-primary supported Brauer classes, not merely primary modulo constants.
It follows that
\begin{equation}\label{eq:regular-supported-local-annihilator}
 j\bigl(C(\mathbb Q_p)^{\operatorname{Br}_{\{p\}}(\mathcal C)[p^\infty]}\bigr)
       \subset B.
\end{equation}
The quotient map extends to proper models away from $p$, so Brauer
functoriality applies to it as in that paper's equation~(3).

Consider the natural map
\[
 \chi:A(\mathbb Q_p)/pA(\mathbb Q_p)
                \longrightarrow A(K_\lambda)/\pi A(K_\lambda).
\]
Its image is the one-dimensional $\Delta$-invariant subspace,
and it is nonzero
\citep[Corollary~2 and Proposition~5, pp.~574--575]{McCallum1994}.
We claim that the saturation $B^{\mathrm{sat}}$ lies in
$\ker\chi$.  If $rP\in B$ for a nonzero integer $r$, set
$t=v_\pi(r)$ and pair $rP$ with $d_{t+1,\lambda}$.
Writing $r=\pi^tu$ and moving the endomorphism across the pairing
reduces the result to a nonzero scalar multiple of the pairing
of $\chi(P)$ with $d_{1,\lambda}$.  Indeed the Rosati dual of
$\pi$ is a unit times $\pi$, and every unit multiplier reduces
modulo $\pi$ to a scalar in $\mathbb F_p$.
Since $P$ is fixed by $\Delta$, its class remains in the invariant
line.  The nontrivial pairing on that line forces $\chi(P)=0$.
This is the algebra in
\citet[Theorem~3, pp.~575--576]{McCallum1994}; equation~\eqref{eq:regular-supported-local-annihilator}
replaces the rational-point reciprocity condition in its proof.

Here $B$ is closed.  After quotienting torsion, $A(\mathbb Q_p)$
is a free $\mathbb Z_p$-module of rank $g=\dim A$, and the
saturation of a closed submodule is again closed.  Since
$B^{\mathrm{sat}}\bmod p\subset\ker\chi$, elementary lattice
linear algebra gives a saturated subgroup $M$ containing
$B^{\mathrm{sat}}$, with $M/A(\mathbb Q_p)_{\mathrm{tors}}$
of rank $g-1$ and image modulo $p$ equal to $\ker\chi$.
For example, extend a basis of the saturated submodule to a
basis of the free module and lift the nonzero functional
$\chi$ on the resulting quotient.

The logarithmic pairing with invariant differentials has torsion
kernel on the left.  Its annihilator of $M$ is a line; choose
a primitive element $\eta$ in the integral N\'eron-differential
lattice on this line and put $\omega=j^*\eta$.
The integral pairing and the identification of N\'eron and
dualizing differential lattices are
\citet[Propositions~2, 3 and~7, pp.~569--570, 578]{McCallum1994}.
Thus the nonzero reduction of $\omega$ belongs to the subspace
$V$ annihilating $\ker\chi$ in that paper.  By its
Theorem~4, p.~579, no nonzero member of $V$ vanishes at all
three branch cusps.

The Coleman integral $\lambda_\omega$, normalized at $O_0$,
vanishes on $j^{-1}(B)$ and hence on the local Brauer zero set.
It vanishes at all three rational branch cusps because their
Abel--Jacobi differences are torsion.  At a branch cusp where
the reduced differential is nonzero, the integral has at most
one zero in that residue disc by
\citet[Lemma~7, pp.~579--580]{McCallum1994}.
Consequently that cusp is the only point of the local Brauer
zero set in its disc.

The quotient $F_p\to C$ is injective on $\mathbb Q_p$-points:
the quotient coordinates determine the $p$th powers of the
Fermat coordinates, and $\mu_p(\mathbb Q_p)=\{1\}$.
The three cusp residue classes map to the three branch cusp
classes, as follows from the same explicit formulas used in
\citet[Lemma~8, p.~580]{McCallum1994}.
Therefore one coordinate-cusp disc on $F_p$ contains just its
cusp in the Brauer zero set.  Rational coordinate permutations,
with the signs appropriate to $X^p+Y^p=Z^p$, preserve the
model away from $p$ and permute the three cusps transitively.
They preserve the local Brauer zero set, proving the first
assertion for all three discs.

For $p=17$, regularity follows directly from Kummer's criterion,
in the full-class-number form of
\citet[Theorem~5.34, p.~78]{Washington1982}.  The relevant
Bernoulli numbers, in increasing even degree, are
\[
 B_2,B_4,B_6,B_8,B_{10},B_{12},B_{14}
 =\frac16,-\frac1{30},\frac1{42},-\frac1{30},
   \frac5{66},-\frac{691}{2730},\frac76.
\]
None has numerator divisible by seventeen.  The wild quotient can be taken to be
$y^{17}=x(1-x)$: here $\gamma=1/4$ and
$4^{16}\equiv154\not\equiv1\pmod {289}$.

It remains to exclude other residue classes.  A primitive
integral point with three unit coordinates would give
$x^{17}+y^{17}=1$ with $x+y\equiv1\pmod {17}$.
Seventeenth powers modulo $289$ depend only on residues
modulo $17$.  For residues $a=0,\ldots,16$, the errors
$a^{17}+((1-a)\bmod17)^{17}-1$ modulo $289$ are, in order,
\[
 \begin{gathered}
 0,0,153,68,102,119,170,34,34,\\
 68,34,34,170,119,102,68,153.
 \end{gathered}
\]
Only $a=0,1$ give zero.  Thus every $\mathbb Q_{17}$-point
reduces to a coordinate cusp.  The first assertion now gives
the stated equality.
\end{proof}

\begin{lemma}[Finite-primary reciprocity]\label{lem:finite-primary-reciprocity}
Let $p$ be an odd prime, let $A/k$ be an abelian variety over a number
field, and let $z\in H^1(k,T_pA)$ have local Kummer representatives
$P_v\in A(k_v)^{\wedge_p}$. For every
$d\in H^1(k,A^\vee)[p^\infty]$,
\[
 \sum_v\langle P_v,d_v\rangle_v=0.
\]
In particular a class supported at one place pairs to zero there.
\end{lemma}
\begin{proof}
Choose $m$ with $p^md=0$. The Kummer sequence supplies a global
lift $b\in H^1(k,A^\vee[p^m])$ of $d$. Reduction of $z$ gives
$z_m\in H^1(k,A[p^m])$. Their Weil cup product lies in
$H^2(k,\mu_{p^m})=\operatorname{Br}(k)[p^m]$.
Its local invariant is, up to one fixed sign convention,
$\langle P_v,d_v\rangle_v$ by the finite Kummer diagram of
\citet[Theorem~I.3.2, Corollary~I.3.4 and Remark~I.3.5,
pp.~41--44]{Milne2006ADT}. Global Brauer reciprocity makes
its finite invariant sum zero; equivalently use that source's
Theorem~I.4.10(b), pp.~56--57, with all places and
$M=A[p^m]$, whose Cartier dual is $A^\vee[p^m]$.
These statements include the residue-characteristic-$p$ place.
The real-place terms vanish because $p$ is odd.
Where $d_v=0$ the pairing is zero even if $b_v$ is a nonzero
Kummer class. This finite-level argument needs no inverse-limit
duality theorem and no finiteness assumption on $\Sha$.
\end{proof}

\begin{corollary}[The primary local cusp condition]
\label{cor:seventeen-primary-local}
Let $z\in H^1(\Q,T_{17}J_{17})$ be locally the Kummer class
of $[x_v-O]$, with $x_v\in X_{17}(\Q_v)$ and a fixed rational
coordinate cusp $O$. Every chosen $x_{17}$ is a rational
coordinate cusp. Consequently this holds for every quotient
Selmer section retaining the compact $T_{17}J_{17}$ class.
The same cusp-disc assertion holds for every regular $p>5$;
if every $\Q_p$-point reduces to the coordinate boundary, the
entire local representative is a rational cusp.
\end{corollary}
\begin{proof}
Push $z$ to the Jacobian of a wild cyclic quotient used in
Proposition~\ref{prop:regular-primary-brauer-cusp-discs}, and
restrict to $K=\Q(\zeta_p)$. Apply
Lemma~\ref{lem:finite-primary-reciprocity} to every $d_k$, using
the rational principal polarization. There is only one place
of $K$ above $p$ and the other local classes vanish. The chosen
local point therefore belongs to exactly the common annihilator
$B$ used in that proposition. Pulling back the classes and
permuting the coordinate cusps gives its cusp-disc conclusion.
At seventeen the displayed residue calculation exhausts all
local points. A quotient section supplies $z$ by filling the
boundary and abelianizing; a full local cusp witness fills to
its center, so every chosen witness has the required Kummer class.
No all-prime descent or global full-section lift is involved.
The older finite-descent consequence also follows, since
\citet[Corollary~7.3, p.~379]{Stoll2007} identifies proper finite
abelian descent with the Brauer--Manin condition.
\end{proof}

\subsection{Nonabelian localization at every odd regular prime}
\label{subsec:regular-nonabelian}

\begin{proposition}[The regular-prime arithmetic source]
\label{prop:regular-nonabelian-localization}
Let $p$ be odd with $p\nmid h(K)$, where $K=\Q(\zeta_p)$,
and let $n=p^r$, $r\ge1$. Put
$N_n=\pi_1(U_{n,\overline\Q})^{\mathrm{pro}\text{-}p}$ and
$C_n=\pi_1(X_{n,\overline\Q})^{\mathrm{pro}\text{-}p}$.
For a standard rational tangent the actual actions on these
and all their stable quotients factor through
$\Gamma=\Gal(K_{\{p\}}(p)/\Q)$, where $K_{\{p\}}(p)$ is the
maximal pro-$p$ extension unramified outside the sole prime above $p$.
Their $G_K$-images are pro-$p$, and $G_{\Q_p}\to\Gamma$ is surjective.
For every finitely generated pro-$p$ group $M$ with such an action,
\[
 H^1(G_{\Q,\{p\}},M)\longrightarrow H^1(\Q_p,M)
\]
is injective as a map of nonabelian cohomology sets.
Localization is consequently injective on proper point-local
Selmer sections for every stable quotient of $C_n$, and on
$H^1(G_{\Q,\{p\}},T_pJ_n)$ and its compact Selmer subgroup.
\end{proposition}
\begin{proof}
In the minus-sign model, put $t=x^n$, $y^n=1-t$. The geometric
pro-$p$ group $\Pi$ of $\mathbb P^1-\{0,1,\infty\}$ at its
unit tangent is free of rank two, with
$\Pi^{\mathrm{ab}}=\Z_p(1)^2$. The coordinate cover has group
$(\Z/n)^2$ and full geometric group
\[
 N_n=\ker(\Pi\longrightarrow\Pi^{\mathrm{ab}}/n\Pi^{\mathrm{ab}}).
\]
Indeed, the core of the kernel of any finite $p$-quotient of the
covering subgroup has $p$-group kernel and $p$-group deck quotient;
thus the induced topology is the full pro-$p$ topology. Riemann
existence supplies this covering comparison
\citep[Expos\'e~XII, Theorem~5.1, p.~251]{SGA1}.
The branch $x=t^{1/n}$, $y=(1-t)^{1/n}$ with leading coefficient
and constant term one gives compatible rational tangential splittings.
The upstairs action is the restriction of the actual action on $\Pi$.

On $\Pi/\Phi(\Pi)$ the action is cyclotomic, so $G_K$ acts
trivially. The automorphism kernel on this Frattini quotient is
pro-$p$: for a finite $p$-group, automorphisms trivial on the
Frattini quotient act freely on the $p$-power set of tuples lifting
a fixed generating basis. Pass to characteristic finite quotients.
Thus the whole $G_K$-image is pro-$p$, also on subgroups and stable
quotients. The based action is unramified outside $p$ by
\citet[Definitions~17.4--17.6, Theorem~17.7 and Corollary~17.8,
pp.~49--51]{Wojtkowiak2009}: the unit tangent extends nontrivially
on the smooth pair there. This is an actual action statement.
Filling the punctures gives the proper action at the rational center.

Let $P=\Gal(K_{\{p\}}(p)/K)$ and let $D$ be the image of the
sole local group $G_{K_\lambda}$. If $D$ were proper, a finite
$p$-quotient and a maximal subgroup would give a cyclic degree-$p$
extension of $K$, unramified outside $\lambda$ and split there.
It would be everywhere unramified, contradicting $p\nmid h(K)$
by class field theory. Hence $D=P$. Total ramification of $p$ in
$K/\Q$ makes $G_{\Q_p}$ surject onto $\Gal(K/\Q)$ as well, proving
surjectivity onto $\Gamma$. The pro-$p$ unramified action over
$K$ factors through $P$, and hence the full action factors through
$\Gamma$.

For any cocycle $c$ on $G_{\Q,\{p\}}$, its restriction to $G_K$
has graph $g\mapsto(c(g),\rho(g))$ in the pro-$p$ semidirect
product $M\rtimes\rho(G_K)$. It factors through $P$. Thus $c$
and its action factor through $\Gamma$. If two classes agree at
$p$, conjugate one cocycle globally by the same local conjugator.
The cocycles now agree on the surjective image of $G_{\Q_p}$,
so agree globally. This proves injection of the entire pointed set.
The identical graph argument applies to finite free $\Z_p$-modules
obtained by subquotients and tensor operations from $N_n^{\mathrm{ab}}$.

For the proper group, smooth proper specialization at $q\ne p$
identifies the pointed arithmetic sequence with the pullback of
the special-fibre sequence
\citep[Expos\'e~X, Theorem~3.8 and Corollary~3.9, p.~217]{SGA1}.
Each proper local point extends and gives a cocycle trivial on
inertia. Filling a local cusp gives its proper point section.
Since the actual coefficient inertia is trivial, conjugacy preserves
this property. Proper Selmer cocycles therefore descend to
$G_{\Q,\{p\}}$. Compare two of them using the same fixed rational
cusp as origin, and apply the proved injection. This also covers
all stable proper quotients. Abelianization is $T_pJ_n$
\citep[Expos\'e~IX, \S5.8, p.~194]{SGA1}, proving the integral
compact assertions, without a rank or $\Sha$ hypothesis.
\end{proof}

\begin{remark}[Why one place does not give the all-prime theorem]
\label{rem:regular-all-prime-limitation}
The primary local condition and nonabelian localization do not imply
$X_{17}(\mathbb A_{\mathbb Q})_\bullet^{f\text{-}\mathrm{ab}}
=X_{17}(\mathbb Q)$.  At positive Mordell--Weil rank, a
single finite place cannot detect the full profinite
Mordell--Weil completion.  Indeed $A(\mathbb Q_q)$ has an
open pro-$q$ subgroup with finite quotient $G$.  For every
prime $\ell\ne q$ not dividing $|G|$, its $\ell$-adic
completion is zero, whereas $A(\mathbb Q)^{\wedge}$ contains
$\mathbb Z_\ell^{\operatorname{rank}A(\mathbb Q)}$.
The same obstruction applies to any fixed finite collection
of local places.

There is also a curve-level distinction.  Given a rational
point $O$ on a smooth proper curve, fix a good prime $q\ne p$
and choose a different $Q_q$ sufficiently close to $O$.
The Abel--Jacobi difference lies in the pro-$q$ formal
subgroup, where multiplication by $p$ is an automorphism.
The adelic point equal to $O$ except for the component $Q_q$
therefore has trivial compact $p$-adic Kummer class at every
place and survives every $p$-primary finite abelian torsor.
It need not be diagonal.  This is a limitation of the
single-primary condition, not a counterexample to full
finite abelian descent.  Stoll's Theorem~8.2 requires finite
support at a density-one set of places; the one place
seventeen does not meet that hypothesis after a finite
number-field extension.
\end{remark}

\begin{proposition}[A sufficient criterion allowing positive rank]
\label{prop:positive-rank-split-prime-criterion}
Let $X/\mathbb Q$ be a smooth proper curve of genus at least
one, with rational base point $O$.  Let $f:X\to A$ be a
morphism to an abelian variety and $Z\subset X$ a finite
$\mathbb Q$-subscheme.  Suppose that there is a finite
Galois extension $L/\mathbb Q$ such that, for all but finitely
many rational primes $\ell$ split completely in $L$,
\[
 T_\ell\Sha(A/\mathbb Q)=0
\]
and
\[
 \{x\in X(\mathbb Q_\ell):
   \log_\ell(f(x)-f(O))\in
   \operatorname{span}_{\mathbb Q_\ell}
                         \log_\ell A(\mathbb Q)\}
                 \subset Z(\mathbb Q_\ell).
\]
Then
$X(\mathbb A_{\mathbb Q})_\bullet^{f\text{-}\mathrm{ab}}
=X(\mathbb Q)\subset Z(\mathbb Q)$.
\end{proposition}
\begin{proof}
The image in $A$ of a finite-descent adelic point has a
compatible compact $\ell$-adic Selmer class.  The inverse-limit
Kummer sequence and the vanishing of $T_\ell\Sha$ identify
that class with an element of $A(\mathbb Q)^{\wedge_\ell}$.
Its logarithm at $\ell$ lies in the stated rational
Mordell--Weil span.  The local hypothesis therefore places
$x_\ell$ in $Z(\mathbb Q_\ell)$ at every such split prime.

After base change to $L$, the degree-one primes have density
one; outside the ramified primes they lie above precisely
the rational primes splitting completely in $L$.
\citet[Theorem~8.2, pp.~383--384]{Stoll2007} makes the
modified adelic point diagonal in $Z(L)$.  Compatibility
with the original $\mathbb Q$-adeles makes this point
$\mathbb Q$-rational: at a split prime, all its conjugate
components equal the same original local point.
Applying the same theorem over $\mathbb Q$ to this rational
singleton recovers the real components as well. Rational points
survive finite abelian descent and $Z$ is a subscheme of $X$, so
the inclusion is exactly the stated equality. Thus the conclusion
concerns the full modified adelic set, proving
Theorem~\ref{thm:positive-rank-split-intro}.
\end{proof}

The local hypothesis of this criterion has not been proved
at a density-one set of split primes for $X_{17}$.
Proposition~\ref{prop:regular-primary-brauer-cusp-discs} determines
the component at seventeen, while
Corollary~\ref{cor:seventeen-no-rankzero-quotient} explains
why a rank-zero abelian target cannot complete the argument.
